%\part{Referencial Teórico}

\chapter[Referencial Teórico]{Referencial Teórico}

Este capítulo apresenta a fundamentação teórica da pesquisa, abordando os Modelos de Linguagem de Grande Escala (LLMs), suas limitações e aplicações no desenvolvimento de software. Também são discutidos os conceitos de vulnerabilidades de código e o uso do \textit{OWASP Top 10} como referência para avaliação de segurança. Em seguida, são apresentados os princípios da Engenharia de Prompt, com destaque para técnicas como \textit{Chain-of-Thought (CoT)} e Geração Aumentada por Recuperação (RAG), além do papel das ferramentas de Análise Estática de Segurança (SAST) no diagnóstico de falhas. Por fim, são discutidos os trabalhos relacionados que fundamentam a proposta de integração iterativa desenvolvida neste estudo.


\section{Modelos de Linguagem de Grande Escala (LLMs)}

O avanço dos Modelos de Linguagem de Grande Escala, do inglês \textit{Large Language Models} (LLMs), está diretamente associado à consolidação da arquitetura Transformer, proposta por \cite{vaswani_attention_2023}. Diferentemente das redes neurais recorrentes, que processam sequências de forma essencialmente ordenada e dependente de estados anteriores, o \textit{Transformer} utiliza mecanismos de autoatenção, para relacionar diferentes posições de uma sequência de entrada simultaneamente, permitindo que o modelo avalie a relevância de cada token em relação aos demais. Em tarefas envolvendo código-fonte, essa capacidade é especialmente importante, pois uma vulnerabilidade pode depender de relações entre variáveis, chamadas de função, fluxos de dados e decisões de controle distribuídas ao longo de diferentes trechos do programa.

A arquitetura \textit{Transformer} é composta, em sua forma original, por blocos codificadores e decodificadores. O codificador transforma a sequência de entrada em representações contextuais, enquanto o decodificador gera uma sequência de saída de maneira autorregressiva, isto é, prevendo cada novo token com base nos tokens já gerados. A partir dessa estrutura, os modelos de linguagem passaram a ser organizados em três famílias principais: modelos apenas codificadores, modelos apenas decodificadores e modelos codificador-decodificador. Essa diferenciação é relevante para a Engenharia de Software porque cada arquitetura apresenta maior adequação a tipos distintos de tarefas, como compreensão de código, detecção de vulnerabilidades, geração de programas e reparo automático \cite{vaswani_attention_2023}.

Antes de serem utilizados como assistentes de programação, os LLMs passam por um ciclo de treinamento composto por diferentes etapas. Inicialmente, grandes volumes de dados são coletados, limpos, deduplicados e transformados em tokens. Em seguida, ocorre o pré-treinamento, geralmente de forma auto-supervisionada, no qual o modelo aprende regularidades estatísticas da linguagem natural e, no caso dos modelos especializados em programação, também de código-fonte \cite{minaee_large_2025}.

O treinamento em grandes bases de código explica por que esses modelos conseguem gerar programas sintaticamente plausíveis, completar funções, explicar trechos e sugerir correções. No entanto, esse aprendizado não equivale a uma compreensão formal da semântica do programa. O modelo aprende padrões recorrentes presentes nos dados de treinamento e os utiliza para prever a continuação mais provável de uma sequência \cite{minaee_large_2025}. Assim, quando solicitado a corrigir uma vulnerabilidade, o LLM não executa, necessariamente, uma prova de segurança nem garante que a lógica de negócio será preservada. Ele produz uma resposta provável com base no contexto recebido, o que torna essencial a presença de validação externa, especialmente em tarefas críticas \cite{sriram_improving_2026}.

Com o aumento de escala desses modelos, a literatura passou a observar habilidades chamadas emergentes, isto é, capacidades que se tornam mais evidentes em modelos maiores e mais bem treinados. Entre elas, destacam-se o seguimento de instruções, o raciocínio em múltiplas etapas e o aprendizado em contexto. O aprendizado em contexto permite que o modelo adapte sua resposta a uma tarefa específica apenas a partir das instruções, exemplos e informações fornecidas no prompt, sem atualizar seus pesos internos \cite{minaee_large_2025}. Essa característica é fundamental para este trabalho, pois técnicas como \textit{Chain-of-Thought} (CoT) e \textit{Retrieval-Augmented Generation} (RAG) exploram justamente a capacidade do modelo de usar contexto adicional durante a inferência. No caso do reparo de vulnerabilidades, isso significa que o modelo pode receber o alerta emitido por uma ferramenta SAST, trechos de código relacionados, exemplos de correções anteriores, regras de segurança e documentação técnica para orientar sua resposta.

% APROFUNDAR ESTUDOS SOBRE CAMADA BIDIRECIONAL
No contexto de código-fonte, os modelos apenas codificadores, ou \textit{encoder-only}, são voltados principalmente para tarefas de compreensão. Como suas camadas de atenção conseguem acessar a sequência de entrada de forma bidirecional, esses modelos são adequados para analisar o código como um todo e produzir uma representação contextual do trecho analisado \cite{jiang_impact_2023}. Exemplos dessa família incluem CodeBERT e GraphCodeBERT, que são frequentemente aplicados em tarefas de detecção de vulnerabilidades, classificação de código vulnerável e identificação de padrões suspeitos. Entretanto, por não serem projetados originalmente para geração autorregressiva de longas sequências, eles tendem a ser menos adequados para produzir diretamente uma correção completa \cite{zhou_large_2024}.

Os modelos apenas decodificadores, ou \textit{decoder-only}, seguem outra lógica. Eles são autorregressivos e foram projetados para prever o próximo token com base no contexto anterior. Essa característica os torna especialmente eficientes em tarefas de geração contínua, como completar funções, escrever código a partir de descrições em linguagem natural e refatorar trechos. No contexto de desenvolvimento de software, esses modelos como GPT, o Codex, o Code Llama se popularizaram justamente por sua capacidade de produzir código funcional a partir de prompts curtos. No entanto, essa mesma natureza generativa também traz riscos: o modelo pode gerar trechos que parecem corretos, mas contêm vulnerabilidades ou não preservam requisitos implícitos do sistema \cite{minaee_large_2025}.

Já os modelos codificador-decodificador, ou \textit{encoder-decoder}, combinam as duas capacidades anteriores. O codificador lê e representa a entrada, enquanto o decodificador gera uma saída condicionada a essa representação. Por isso, essa arquitetura é especialmente adequada para tarefas de transformação de sequência, nas quais uma entrada deve ser convertida em uma saída relacionada. No campo de código, exemplos importantes incluem CodeT5 e PLBART. Esses modelos podem ser utilizados tanto para compreensão quanto para geração, mas se destacam em tarefas como tradução entre linguagens, sumarização de código, refinamento e reparo de defeitos \cite{minaee_large_2025}.

%ADICIONAR REFERENCIAS DO ZOTERO
%EXPLICAR MELHOR SOBRE A ESCOLHA DO ENCODER-ONLY?
A literatura sobre detecção e reparo de vulnerabilidades confirma essa divisão funcional entre arquiteturas. Estudos de revisão indicam que modelos \textit{encoder-only} aparecem com frequência em tarefas de detecção de vulnerabilidades, enquanto modelos \textit{decoder-only} e \textit{encoder-decoder} são mais recorrentes em tarefas de reparo \cite{zhou_large_2024}. Essa diferença ocorre porque detectar uma vulnerabilidade exige principalmente compreender, representar e classificar um trecho de código, enquanto reparar exige gerar uma nova versão do programa. Embora os modelos \textit{encoder-decoder} pareçam, em princípio, a opção mais completa por combinarem leitura contextual da entrada e geração da saída, eles não são necessariamente a melhor escolha em todos os cenários. Seu desempenho depende de fatores como escala do modelo, dados de pré-treinamento, tipo de ajuste fino, tamanho da janela de contexto e alinhamento com instruções. Além disso, quando tratados como uma arquitetura única para múltiplas tarefas, esses modelos podem apresentar desempenho abaixo do esperado em algumas aplicações específicas. Em tarefas complexas de reparo de vulnerabilidades, modelos \textit{decoder-only} de maior escala e mais bem alinhados para seguir instruções podem superar modelos \textit{encoder-decoder} menores, especialmente quando recebem contexto adicional, exemplos, alertas de ferramentas SAST e feedback iterativo. Portanto, a escolha da arquitetura influencia diretamente o desenho de um \textit{pipeline} de segurança: modelos \textit{encoder-only} tendem a ser mais adequados para triagem e classificação de alertas, modelos \textit{encoder-decoder} são úteis para transformações localizadas de código, e modelos \textit{decoder-only} comerciais ou especializados em código podem ser mais apropriados para a geração de correções em fluxos iterativos e contextualizados. Este trabalho segue essa lógica ao adotar um modelo decoder-only para a etapa de geração das correções, enquanto a triagem e a identificação inicial das vulnerabilidades ficam a cargo da ferramenta SAST, e não de um modelo encoder-only dedicado a essa função.

Apesar desses avanços, o uso de LLMs em tarefas de segurança ainda apresenta limitações relevantes. Um dos problemas centrais é que os modelos aprendem a partir de grandes quantidades de código produzido por desenvolvedores reais, incluindo exemplos corretos, inseguros, desatualizados ou incompletos \cite{chen_evaluating_2021}. Assim, quando não recebem restrições explícitas, podem replicar padrões inseguros presentes nos dados de treinamento. Em vez de selecionar automaticamente a alternativa mais segura, o modelo pode escolher a solução mais comum ou mais provável nos dados em que foi treinado.

Um limite importante aparece quando o modelo é utilizado em uma abordagem \textit{single-pass}, isto é, quando recebe um único prompt e gera uma resposta sem iteração, validação ou realimentação por ferramentas externas. Embora esse tipo de interação seja simples, ela é insuficiente para tarefas de segurança. O reparo de uma vulnerabilidade não consiste apenas em remover o trecho apontado como inseguro, mas em produzir uma modificação que elimine o risco, preserve a lógica original, mantenha a funcionalidade esperada e não introduza novas falhas. Estudos recentes mostram que técnicas de geração segura podem, em alguns casos, aumentar artificialmente a pontuação de segurança ao remover linhas vulneráveis ou produzir código residual sem relação com a tarefa original, fenômeno descrito como \textit{garbage code} \cite{dai_rethinking_2025}. Nesses casos, o código pode parecer seguro para uma análise superficial, mas deixa de cumprir sua finalidade.

Dessa forma, a principal limitação dos LLMs em reparo de vulnerabilidades não está apenas na capacidade de gerar código, mas na ausência de garantias sobre segurança e funcionalidade quando o modelo atua isoladamente. Mesmo modelos avançados podem produzir respostas plausíveis, mas incorretas, inseguras ou incompletas. Por isso, trabalhos recentes defendem o uso de fluxos iterativos com validação externa que forneçam sinais objetivos para orientar a correção. \cite{sriram_improving_2026}, por exemplo, propõem um fluxo com RAG e realimentação por múltiplas ferramentas, incluindo compilador, CodeQL e execução simbólica, reduzindo significativamente a taxa de vulnerabilidades nos programas avaliados. Esse tipo de resultado reforça a premissa deste trabalho: LLMs podem apoiar o reparo automático de vulnerabilidades, mas tendem a ser mais confiáveis quando integrados a ferramentas de análise, enriquecimento de contexto e mecanismos iterativos de validação.

\section{Engenharia de Prompt}

A engenharia de prompt consiste no processo sistemático de estruturar as entradas fornecidas a um Modelo de Linguagem de Grande Escala para induzir comportamentos adequados a uma determinada tarefa. Diferentemente do ajuste fino (\textit{fine tuning}), no qual os parâmetros internos do modelo são modificados a partir de novos exemplos de treinamento, a engenharia de prompt procura explorar os conhecimentos e as capacidades já adquiridos pelo modelo durante seu pré-treinamento. Dessa forma, o modelo pode ser adaptado a diferentes atividades sem a necessidade de realizar um novo treinamento, o que reduz os custos computacionais e permite a utilização de modelos comerciais cujos parâmetros internos não estão disponíveis \cite{reynolds_prompt_2021}.

Um prompt não precisa ser composto apenas por uma pergunta ou ordem direta. Em aplicações mais complexas, ele pode ser estruturado a partir de diferentes componentes, como a definição do papel atribuído ao modelo, a descrição da tarefa, o contexto necessário para executá-la, os dados de entrada, as restrições que devem ser respeitadas, os exemplos de respostas esperadas e o formato de saída \cite{reynolds_prompt_2021}. Em uma tarefa de reparo de vulnerabilidades, por exemplo, o prompt pode informar que o modelo deve atuar como especialista em segurança, apresentar o código analisado, incluir o alerta produzido por uma ferramenta SAST, indicar a CWE relacionada, restringir alterações que não sejam necessárias e exigir que a resposta seja apresentada como um patch ou trecho de código estruturado \cite{khan_code_2025}.

A combinação desses elementos reduz a ambiguidade da tarefa e facilita a integração da resposta em um pipeline automatizado. A definição do formato de saída é especialmente importante porque respostas em linguagem natural, embora compreensíveis para uma pessoa, podem dificultar o processamento automático \cite{liu_we_2024}. Assim, o modelo pode ser orientado a fornecer campos específicos, como a classificação do alerta, a causa da vulnerabilidade, a correção proposta, as linhas modificadas, as justificativas técnicas e eventuais limitações identificadas. A padronização também permite comparar as respostas de diferentes modelos e verificar se todos receberam instruções equivalentes.

Nesse contexto, a engenharia de prompt deve ser compreendida como um processo iterativo e experimental. Pequenas alterações na formulação, na ordem das informações e nos exemplos fornecidos podem modificar significativamente a resposta do modelo \cite{lu_fantastically_2022}.

Convém ainda distinguir engenharia de prompt de \textit{prompt tuning}. A engenharia de prompt utiliza instruções discretas e legíveis escritas em linguagem natural. O termo \textit{prompt tuning}, em seu sentido técnico, costuma representar o aprendizado de vetores contínuos, denominados \textit{soft prompts}, que são ajustados por meio de retropropagação enquanto os parâmetros principais do modelo permanecem congelados \cite{lester_power_2021}. \citeonline{khan_code_2025} utilizam a expressão \textit{context-aware prompt tuning} para descrever o enriquecimento progressivo de prompts textuais com conhecimento de segurança e de código.

Uma das formas mais simples de organizar a interação com uma LLM é o \textit{zero-shot prompting}. Nesse cenário, o modelo recebe a descrição da tarefa, mas nenhum exemplo de como ela deve ser resolvida. Para o reparo de vulnerabilidades, uma instrução zero-shot poderia apenas solicitar que o modelo identifique uma falha no código e proponha uma versão corrigida \cite{reynolds_prompt_2021}. O resultado dependeria integralmente do conhecimento adquirido durante o treinamento e da capacidade do modelo de interpretar a instrução fornecida.

No \textit{one-shot prompting}, é incluído um único exemplo resolvido, enquanto no \textit{few-shot prompting} são apresentados alguns exemplos de entrada e saída antes do problema principal. Esses exemplos não modificam os parâmetros internos do modelo. Eles atuam durante a inferência, ajudando a localizar o comportamento desejado entre os padrões já aprendidos e indicando a estrutura esperada da resposta \cite{reynolds_prompt_2021}. Em um cenário de segurança, um exemplo pode apresentar um alerta SAST, um trecho vulnerável e sua correção correspondente, permitindo que o modelo reconheça como os dados do analisador devem ser interpretados.

Entretanto, a utilização de exemplos não garante automaticamente uma resposta melhor. Exemplos pouco representativos podem induzir o modelo a aplicar uma correção inadequada, copiar padrões irrelevantes ou interpretar detalhes acidentais como parte da tarefa \cite{reynolds_prompt_2021}. Além disso, a literatura mostra que a ordem em que exemplos \textit{few-shot} são apresentados pode alterar substancialmente o desempenho do modelo \cite{lu_fantastically_2022}.

Para tarefas que exigem múltiplas etapas de análise, foi proposta a técnica de Cadeia de Pensamento, ou \textit{Chain-of-Thought (CoT)}. Em vez de solicitar diretamente uma resposta final, o prompt orienta o modelo a decompor o problema em passos intermediários. A formulação original do CoT utiliza exemplos que contêm tanto a entrada quanto uma sequência de etapas que conduz à resposta \cite{wei_chain--thought_2023}. Posteriormente, o chamado zero-shot CoT demonstrou que instruções simples de decomposição, como solicitar que o modelo resolva o problema passo a passo, também podem melhorar o desempenho em determinadas tarefas \cite{kojima_large_2023}.

No reparo de vulnerabilidades, a decomposição pode orientar o modelo a executar uma sequência semelhante à utilizada por um especialista: identificar o trecho apontado pelo SAST, determinar se o alerta representa uma vulnerabilidade real, localizar a origem e o destino do fluxo de dados, identificar a causa da falha, selecionar uma estratégia de mitigação, gerar a correção e verificar se a modificação preserva a funcionalidade original \cite{wei_chain--thought_2023, khan_code_2025}. Essa estrutura é mais adequada do que uma solicitação genérica para “corrigir o código”, pois impede que etapas importantes sejam omitidas e facilita a identificação do ponto em que uma resposta incorreta foi produzida.

Contudo, a Cadeia de Pensamento não deve ser interpretada como uma representação necessariamente fiel do processo interno utilizado pelo modelo. A sequência textual pode constituir uma justificativa plausível produzida após a escolha da resposta, sem refletir exatamente os fatores que determinaram a geração. Portanto, o CoT deve ser utilizado como uma estratégia para estruturar a tarefa e produzir etapas verificáveis, e não como uma garantia de transparência ou correção do raciocínio \cite{lester_power_2021}. Neste trabalho, as propostas de alteração foram avaliadas por reanálise estática e testes automatizados, sem validar o raciocínio interno do modelo.

Uma extensão do CoT é a técnica de \textit{self-consistency}. Em vez de utilizar apenas uma sequência de raciocínio gerada pelo modelo, são produzidas múltiplas respostas para o mesmo problema. Em seguida, seleciona-se a resposta final que apresenta maior concordância entre os diferentes caminhos gerados. \citeonline{wang_self-consistency_2023} demonstraram que essa estratégia pode melhorar significativamente o desempenho em tarefas de raciocínio, sem exigir treinamento ou ajuste fino adicional.

A aplicação direta de \textit{self-consistency} ao reparo de código, entretanto, exige adaptações. Uma vulnerabilidade pode admitir várias correções textualmente diferentes, mas semanticamente equivalentes. Assim, não seria adequado selecionar uma solução apenas pela igualdade entre os códigos gerados \cite{wang_self-consistency_2023}. Uma possível aplicação consiste em gerar múltiplos candidatos e utilizar critérios externos para compará-los, como compilação, execução de testes, quantidade de alertas SAST restantes, ausência de novas vulnerabilidades e preservação do comportamento esperado. Nesse cenário, a ferramenta de análise e os testes assumem o papel de mecanismos de seleção e validação das alternativas.

Outra estratégia relevante para fornecer informações ao modelo é a Geração Aumentada por Recuperação, ou \textit{Retrieval-Augmented Generation} (RAG). Diferentemente de uma técnica composta apenas pela escrita de instruções, o RAG constitui uma arquitetura que conecta a LLM a fontes externas de conhecimento. Sua proposta original combina a memória paramétrica, correspondente às informações representadas nos parâmetros do modelo, com uma memória não paramétrica, composta por documentos armazenados externamente e recuperados no momento da consulta \cite{klesel_retrieval-augmented_2025}.

Um pipeline de RAG pode ser dividido em uma etapa de preparação e outra de consulta. Na preparação, os documentos são coletados, processados e separados em fragmentos, frequentemente denominados \textit{chunks}. Cada fragmento é convertido em uma representação vetorial por um modelo de \textit{embeddings} e armazenado em uma base vetorial. Durante a consulta, a entrada do usuário também é transformada em um vetor, permitindo que o recuperador calcule a similaridade entre a pergunta e os documentos armazenados. Os fragmentos mais relevantes são selecionados, incorporados ao prompt e enviados à LLM junto com a tarefa original. Por fim, o modelo gera uma resposta condicionada tanto ao seu conhecimento interno quanto ao conteúdo recuperado \cite{klesel_retrieval-augmented_2025}.

As três etapas centrais desse processo são a recuperação, o enriquecimento e a geração. Na recuperação, são localizados os documentos considerados mais relevantes para a consulta. No enriquecimento, os documentos selecionados são organizados e anexados à instrução. Na geração, a LLM produz a resposta utilizando o contexto ampliado \cite{klesel_retrieval-augmented_2025}. Neste trabalho, a recuperação utilizou um mapeamento entre CWEs e documentos de apoio, com alternativas por categoria e documento genérico, conforme descrito na metodologia.

A base de conhecimento pode incluir descrições do CWE, recomendações do OWASP, regras e documentações do SonarQube, guias de programação segura, documentação oficial da linguagem e do framework analisado, além de exemplos de reparos previamente validados. Dessa forma, quando uma ferramenta SAST identificar uma falha criptográfica, por exemplo, o sistema poderá recuperar informações sobre algoritmos recomendados, APIs atuais e formas adequadas de armazenamento de chaves. Para uma vulnerabilidade de injeção, poderão ser recuperadas recomendações relacionadas a consultas parametrizadas, validação de entrada e neutralização de dados.

Embora o RAG possa melhorar a fundamentação das respostas e fornecer informações que não estavam presentes no treinamento do modelo, ele não elimina automaticamente as alucinações. O resultado depende da qualidade da base documental, da divisão dos documentos, do modelo de \textit{embeddings}, da formulação da consulta e da capacidade do recuperador de selecionar informações realmente relevantes. Caso um documento incorreto, desatualizado ou pouco relacionado seja recuperado, ele poderá induzir a LLM a gerar uma correção inadequada \cite{klesel_retrieval-augmented_2025}. Portanto, os documentos utilizados devem ser provenientes de fontes confiáveis e sua origem deve ser mantida para permitir rastreabilidade e verificação.

No domínio de segurança, a importância do contexto fornecido pelo prompt é demonstrada por \cite{khan_code_2025}. Ao avaliarem o GitHub Copilot em 156 trechos com vulnerabilidades relacionadas a estouro de buffer, os autores observaram uma taxa de 76\% de sucesso na detecção, mas apenas 15\% no reparo quando nenhum conhecimento adicional era fornecido. Informar somente que o código possuía uma vulnerabilidade elevou a taxa de reparo para 20\%, enquanto a inclusão do tipo de CWE aumentou o resultado para 31\%. Quando o prompt passou a incorporar progressivamente informações sobre segurança e sobre o contexto do código, a taxa de sucesso alcançou 63\%.

Neste trabalho, a engenharia de prompt organizou a instrução, o código e as restrições de alteração. Conforme a estratégia, foram acrescentados um roteiro de raciocínio guiado, o diagnóstico SAST e um documento recuperado pelo RAG simplificado. Essas combinações constituíram as estratégias P1--P4, descritas na metodologia.

As respostas foram solicitadas em formato JSON, contendo as edições propostas, suas justificativas e limitações. As propostas aplicadas foram avaliadas por compilação, testes automatizados e reanálise estática. As justificativas produzidas pela LLM não foram utilizadas como evidência suficiente de validade da correção.

O enriquecimento de contexto e o raciocínio guiado apoiaram a geração das propostas, enquanto a análise estática e os testes determinaram sua avaliação e a continuidade da sequência. Cada tentativa partiu do código original, sem incorporar ao prompt as respostas ou os erros das tentativas anteriores. A iteração consistiu, portanto, na execução de estratégias sucessivas até atender à regra de parada ou esgotar P1--P4.

\section{Vulnerabilidades de Software e OWASP 10}

% ADICIONAR REFERÊNCIAS DO ZOTERO
Uma vulnerabilidade de software pode ser compreendida como uma fraqueza presente na arquitetura, no projeto, na implementação ou na configuração de um sistema que pode ser explorada por um agente malicioso para violar alguma propriedade de segurança. A exploração dessas falhas pode comprometer a confidencialidade, a integridade ou a disponibilidade das informações e dos serviços oferecidos pelo sistema. Dependendo da vulnerabilidade, as consequências podem incluir acesso indevido a dados, execução não autorizada de código, alteração de informações e interrupção de serviços \cite{calin_software_2020}.

Entre as consequências mais relevantes encontram-se a negação de serviço, do inglês \textit{Denial of Service} (DoS), e a escalação de privilégios. Em um ataque de negação de serviço, uma falha é explorada para tornar determinado recurso ou serviço indisponível aos usuários legítimos, por exemplo, provocando o encerramento da aplicação ou consumindo excessivamente seus recursos. Já na escalação de privilégios, o atacante obtém permissões superiores às que deveria possuir, podendo passar de um usuário comum para um usuário administrativo e, assim, ler, modificar ou excluir informações protegidas. \citeonline{calin_software_2020} identificaram essas consequências entre as mais relevantes nas vulnerabilidades de maior impacto analisadas em seu estudo.

Para organizar a identificação e a comunicação dessas falhas, a área de segurança utiliza um ecossistema composto por diferentes bases e sistemas de pontuação. Entre os principais componentes desse ecossistema encontram-se o \textit{Common Vulnerabilities and Exposures} (CVE), a \textit{National Vulnerability Database} (NVD), o \textit{Common Weakness Enumeration} (CWE) e o \textit{Common Vulnerability Scoring System} (CVSS). Embora esses mecanismos sejam relacionados, cada um desempenha uma função diferente no gerenciamento de vulnerabilidades.

O ecossistema de gerenciamento de vulnerabilidades é composto por mecanismos complementares. O Common Vulnerabilities and Exposures (CVE) identifica vulnerabilidades publicamente divulgadas por meio de códigos únicos, permitindo que pesquisadores, fabricantes e ferramentas se refiram à mesma falha \cite{cve_overview}. A National Vulnerability Database (NVD), mantida pelo NIST, reúne esses registros e os enriquece com informações como produtos afetados, referências, configurações vulneráveis e métricas de severidade \cite{nvd_home}. Já o Common Weakness Enumeration (CWE) classifica os tipos de fraqueza que podem originar vulnerabilidades, permitindo que diferentes CVEs sejam associados a uma mesma causa estrutural \cite{cwe_about}.

O Common Vulnerability Scoring System (CVSS) complementa esse ecossistema ao atribuir uma pontuação de severidade entre 0 e 10, classificada como nenhuma, baixa, média, alta ou crítica. Assim, enquanto o CVE identifica uma vulnerabilidade concreta, o CWE descreve o tipo de fraqueza, a NVD agrega informações sobre o registro e o CVSS comunica sua severidade técnica. Entretanto, essa pontuação não representa, isoladamente, o risco para uma organização, que também depende da exposição do componente, da importância do ativo, da possibilidade de exploração, dos controles existentes e do impacto para o negócio \cite{first_cvss_v4}.

No contexto específico de aplicações web, uma das principais referências para conscientização e orientação é o OWASP Top 10. Mantido pela Open Worldwide Application Security Project (OWASP), o documento apresenta categorias que representam um consenso amplo sobre os riscos mais críticos enfrentados por aplicações web. O OWASP Top 10 não deve ser interpretado como uma lista completa de todas as vulnerabilidades existentes nem como uma norma de certificação. Seu objetivo principal é ampliar a conscientização, apoiar treinamentos e orientar organizações na identificação e prevenção de classes recorrentes de problemas de segurança \cite{owasp_top10_2025_intro}.

A construção da edição de 2025 foi fundamentada em uma combinação de dados empíricos e avaliação da comunidade especializada. Organizações de segurança forneceram dados provenientes de testes realizados em mais de 2,8 milhões de aplicações. Diferentemente da edição de 2021, a coleta de 2025 não restringiu previamente quais CWEs poderiam ser enviadas, resultando na análise de 589 tipos de fraquezas. Para cada CWE, o OWASP considerou a quantidade de aplicações testadas e a quantidade de aplicações em que pelo menos uma ocorrência da fraqueza foi identificada \cite{owasp_top10_2025_intro}.

Essa metodologia considera principalmente a prevalência da fraqueza entre as aplicações, e não o número absoluto de ocorrências dentro de uma mesma aplicação. Dessa forma, um sistema contendo quatro instâncias de determinada CWE e outro contendo quatro mil instâncias são contabilizados igualmente para a medição de prevalência. Essa decisão reduz distorções entre testes manuais, que frequentemente registram uma falha apenas uma vez, e ferramentas automatizadas, que podem registrar separadamente cada ocorrência encontrada \cite{owasp_top10_2025_intro}. 

Além dos dados quantitativos, o OWASP utiliza uma pesquisa com profissionais de desenvolvimento e segurança. A edição de 2025 é descrita como informada por dados, mas não cegamente dirigida por eles. Isso ocorre porque os resultados fornecidos pelas ferramentas representam, em grande parte, vulnerabilidades que a indústria já sabe detectar em escala. Novos riscos podem demorar anos até serem incorporados aos mecanismos automatizados de teste \cite{owasp_top10_2025_intro}.

Outro aspecto relevante da edição de 2025 é a priorização das causas raízes das vulnerabilidades em vez de seus sintomas. Um sintoma representa o efeito observável de uma falha, enquanto a causa raiz corresponde à condição técnica que permitiu sua ocorrência \cite{owasp_top10_2025_intro}. Essa escolha possui importância prática para o reparo automático. Uma descrição baseada somente no sintoma pode indicar que dados foram expostos, mas fornece poucas informações sobre qual modificação deve ser realizada no código. Já a identificação da causa raiz pode indicar a necessidade de substituir um algoritmo obsoleto, remover uma chave inserida diretamente no código ou utilizar uma função segura para armazenamento de senhas, por exemplo. Dessa maneira, categorias orientadas à causa oferecem informações mais adequadas para ferramentas SAST e modelos de linguagem responsáveis por sugerir correções \cite{khan_code_2025}. 

Entre as categorias que compõem o OWASP Top 10 de 2025, a Quebra de Controle de Acesso, as Falhas Criptográficas, as vulnerabilidades de Injeção, as Falhas de Autenticação e o Tratamento Incorreto de Condições Excepcionais merecem atenção por afetarem mecanismos essenciais de segurança das aplicações e por poderem comprometer propriedades como confidencialidade, integridade, disponibilidade e autenticidade. Essas categorias representam diferentes formas de fragilidade no desenvolvimento de software e são apresentadas a seguir com maior detalhamento quanto às suas características, formas de ocorrência e possíveis impactos \cite{owasp_a01_2025,owasp_a04_2025,owasp_a05_2025,owasp_a07_2025,owasp_a10_2025}.

 

A A01:2025 (Quebra de Controle de Acesso) ocorre quando as restrições sobre as ações permitidas aos usuários não são aplicadas corretamente, possibilitando o acesso a recursos de terceiros, páginas administrativas, manipulação de tokens e elevação indevida de privilégios; essa categoria também passou a incluir ataques de\textit{ Server-Side Request Forgery }(SSRF). 

A A04:2025 (Falhas Criptográficas) abrange problemas como ausência de criptografia, uso de algoritmos ou protocolos obsoletos, exposição de chaves, armazenamento inadequado de senhas e geração de valores previsíveis, podendo resultar na exposição de dados sensíveis e no comprometimento do sistema. 

A A05:2025 (Injeção) ocorre quando entradas não confiáveis são enviadas a interpretadores sem validação, neutralização ou parametrização adequadas, incluindo ataques como \textit{SQL Injection, NoSQL Injection}, injeção de comandos e \textit{Cross-site Scripting} (XSS); sua prevenção envolve consultas parametrizadas, validação de entradas e neutralização contextual dos dados. 

A A07:2025 (Falhas de Autenticação) envolve o reconhecimento indevido de usuários, a proteção inadequada de credenciais e sessões, o uso de senhas incorporadas ao código, credenciais padrão e mecanismos vulneráveis à fixação de sessão, sendo recomendada a utilização de \textit{frameworks} consolidados e autenticação multifator. 

Por fim, a A10:2025 (Tratamento Incorreto de Condições Excepcionais) está relacionada a erros não tratados, exceções não capturadas, falhas lógicas e situações em que o sistema falha de maneira insegura, podendo provocar vazamento de informações, corrupção do estado da aplicação, transações inconsistentes ou negação de serviço \cite{owasp_top10_2025_intro}.

Essa diversidade de vulnerabilidades possibilita analisar não apenas se a LLM remove o alerta produzido pela ferramenta SAST, mas também se ela compreende a causa da vulnerabilidade e gera uma correção que mantém a funcionalidade originalmente esperada.

\section{Ferramentas de Análise Estática de Segurança (SAST)}

O Teste Estático de Segurança de Aplicações, do inglês \textit{Static Application Security Testing} (SAST), compreende um conjunto de técnicas e ferramentas utilizadas para examinar o código-fonte de um sistema sem executar a aplicação. Seu objetivo é identificar e relatar fraquezas de implementação que possam originar vulnerabilidades, como manipulação insegura de entradas, uso inadequado de memória, chamadas perigosas de funções, configurações criptográficas frágeis e violações de práticas de programação segura. Por ter acesso direto ao código, essa abordagem também é classificada como uma forma de teste de caixa-branca \cite{noauthor_static_nodate,nist_source_code_security_analyzers_2021}.

A análise estática é normalmente aplicada durante a fase de implementação do ciclo de desenvolvimento seguro de software. Como não exige que o sistema esteja implantado e em execução, ela pode ser executada desde as primeiras versões do código e repetida a cada alteração realizada pelos desenvolvedores. Essa característica favorece a adoção do princípio de \textit{shift-left}, no qual atividades de segurança são antecipadas para etapas iniciais do desenvolvimento, reduzindo o custo de correções que, de outra forma, seriam descobertas apenas durante testes finais ou após a implantação do sistema \cite{noauthor_static_nodate}.

As ferramentas SAST também podem ser integradas a ambientes de desenvolvimento, repositórios de código e pipelines de integração contínua. Dessa forma, um novo trecho pode ser analisado durante sua escrita, em uma solicitação de integração ou antes da publicação de uma nova versão. A repetibilidade representa uma das principais vantagens da análise estática, pois permite examinar grandes quantidades de código de forma automática e executar verificações frequentes, como em cada \textit{commit} e \textit{merge request}. Entretanto, essa automação não significa que todos os alertas correspondam necessariamente a vulnerabilidades reais \cite{iranmanesh_zerofalse_2025}.

O funcionamento de uma ferramenta SAST pode envolver diferentes técnicas derivadas de compiladores e de análise de programas. Essas técnicas são utilizadas em conjunto para representar a estrutura do código, acompanhar as possíveis formas de execução e verificar se determinados padrões violam regras de segurança. Entre as técnicas mais relevantes estão a análise léxica, a análise de fluxo de controle, a análise de fluxo de dados e a análise de propagação de dados não confiáveis, conhecida como \textit{taint analysis} \cite{noauthor_static_nodate}.

A análise léxica transforma o código-fonte em unidades menores chamadas tokens. Palavras-chave, identificadores, operadores, literais e símbolos são separados e classificados, permitindo que a ferramenta deixe de tratar o programa como uma sequência comum de caracteres e passe a reconhecer sua estrutura sintática. Essa etapa permite detectar padrões locais, como a utilização de funções consideradas perigosas, credenciais inseridas diretamente no código, algoritmos criptográficos obsoletos ou construções incompatíveis com determinadas regras de programação segura \cite{noauthor_static_nodate}. Neste trabalho, a análise léxica é realizada inteiramente pela ferramenta SAST como etapa interna de seu próprio mecanismo de detecção.

Entretanto, vulnerabilidades complexas não podem ser detectadas apenas pela procura de palavras ou padrões isolados. Para compreender os diferentes caminhos que um programa pode seguir, as ferramentas podem construir uma representação de seu fluxo de controle. Nessa representação, o código é dividido em blocos básicos, formados por sequências de instruções executadas consecutivamente, e as transições possíveis entre esses blocos são organizadas em um Grafo de Fluxo de Controle, do inglês Control Flow Graph (CFG). Esse grafo permite representar condições, repetições, chamadas e desvios que influenciam a execução do programa \cite{noauthor_static_nodate}.

A análise de fluxo de dados permite que ferramentas SAST acompanhem como valores são definidos, modificados e utilizados ao longo dos possíveis caminhos do programa. Uma aplicação importante dessa técnica é a \textit{taint analysis}, que identifica dados provenientes de fontes não confiáveis e rastreia sua propagação até operações sensíveis, denominadas \textit{sinks}. Caso esses dados alcancem um destino, como uma consulta SQL, uma execução de comando ou a geração de uma página web, sem passar por validação, sanitização ou neutralização adequada, a ferramenta pode emitir um alerta de vulnerabilidade. Dessa forma, a análise considera a relação entre fonte, sanitizador e destino, permitindo identificar não apenas a linha problemática, mas todo o caminho percorrido pela entrada. Essas técnicas possibilitam detectar antecipadamente vulnerabilidades como injeção SQL, \textit{Cross-site Scripting} e uso de funções inseguras, sem exigir a execução da aplicação ou a preparação de entradas específicas de teste \cite{noauthor_static_nodate}.

Apesar dessas vantagens, a análise estática apresenta limitações. Um resultado produzido pela ferramenta representa a identificação de uma condição potencialmente problemática, e não necessariamente a comprovação de que ela poderá ser explorada em uma aplicação real. Um falso positivo ocorre quando a ferramenta aponta uma possível vulnerabilidade que, após análise, não representa um risco real. Isso pode acontecer, por exemplo, quando a ferramenta não reconhece um mecanismo de validação personalizado utilizado pela aplicação ou não consegue acompanhar o comportamento de um componente externo \cite{iranmanesh_zerofalse_2025}.

O ecossistema de análise estática inclui ferramentas comerciais, acadêmicas e de código aberto, com diferentes níveis de especialização. Algumas possuem cobertura para diversas linguagens e tipos de fraqueza, enquanto outras são direcionadas a uma única linguagem, plataforma ou categoria de vulnerabilidade \cite{nist_source_code_security_analyzers_2021}. A escolha de uma ferramenta deve considerar fatores como as linguagens suportadas, os tipos de vulnerabilidade detectados, a necessidade de o projeto estar completamente compilável, a integração com IDEs e pipelines, os formatos de relatório disponíveis e a possibilidade de personalizar regras.

Em ferramentas de análise estática, uma regra representa uma condição ou prática que deve ser verificada no código. Quando o analisador encontra uma violação, é criado um \textit{issue} associado à regra responsável pela detecção. As regras podem ser organizadas de acordo com aspectos de segurança, confiabilidade e manutenibilidade \cite{noauthor_sonarqube_2026}.

As regras possuem ainda estados que indicam seu grau de maturidade. Uma regra classificada como \textit{Ready} é considerada adequada para utilização em produção; uma regra Beta ainda não recebeu feedback suficiente e pode apresentar falsos positivos ou negativos; e uma regra \textit{Deprecated} não deve mais ser utilizada quando existe uma alternativa mais atual ou precisa. As regras também podem ser associadas a níveis de severidade, definidos a partir do impacto potencial e da probabilidade de ocorrência ou exploração do problema \cite{noauthor_sonarqube_2026}.

No pipeline implementado, a análise inicial com Semgrep estabeleceu os alertas de referência e forneceu os dados utilizados na constituição do catálogo. Nas estratégias P3 e P4, o diagnóstico do alerta também integrou o contexto enviado à LLM. Após a aplicação de cada proposta, a reanálise verificou a permanência do alerta original e o aparecimento de novos alertas. Essa avaliação foi combinada com compilação e testes automatizados, pois o desaparecimento de um alerta não comprova, isoladamente, uma correção segura e funcional.

Nota-se, portanto, que a principal contribuição do SAST para o pipeline não se limita a informar que existe um possível problema. A ferramenta fornece um diagnóstico estruturado que pode reduzir a ambiguidade da tarefa de reparo. Informações sobre a regra violada, o tipo de fraqueza, a localização e o fluxo de dados ajudam a LLM a identificar a causa da vulnerabilidade e a escolher uma correção compatível com o código analisado.


\section{Trabalhos Relacionados}

A literatura recente tem demonstrado que o uso isolado de Modelos de Linguagem de Grande Escala (abordagem \textit{single-pass}) é insuficiente para a correção confiável de vulnerabilidades de software. Estudos apontam que, sem o fornecimento de um contexto adequado, os modelos falham em compreender a causa raiz do problema de segurança e frequentemente comprometem a funcionalidade do sistema ao tentar corrigi-lo \cite{dai_rethinking_2025}.

Para otimizar o controle e a eficácia dessas ferramentas, \citeonline{reynolds_prompt_2021} introduziram o conceito formal de programação de prompts demonstrando que a forma como a instrução é estruturada em linguagem natural impacta drasticamente o desempenho do modelo. Os autores ressaltam a importância de técnicas nas quais a inteligência artificial é induzida a decompor um problema complexo em etapas menores antes de um veredito, evitando respostas precipitadas e permitindo um processamento lógico mais eficiente.

Diante deste contexto de reparos, novas metodologias passaram a integrar técnicas de Geração Aumentada por Recuperação (RAG) e Cadeia de Pensamento (CoT). O estudo de \citeonline{tao_adapting_2025} combinou essas estratégias para orientar o modelo a primeiro compreender a semântica do código para, em seguida, analisar os trechos vulneráveis e gerar sugestões de reparo. A pesquisa demonstrou que integrar a recuperação de conhecimentos externos (RAG) juntamente ao raciocínio em etapas (CoT) alcança um desempenho superior em diversas métricas na geração de instruções e soluções de segurança.

Paralelamente à evolução da engenharia de prompts, os estudos evidenciam que os modelos generativos atingem sua máxima eficácia quando inseridos em pipelines iterativos com ferramentas de análise externa. \citeonline{sriram_improving_2026} desenvolveram um fluxo de auto-reparo no qual um LLM refina continuamente suas próprias respostas utilizando o feedback diagnóstico de ferramentas SAST (como o CodeQL). Esse estudo comprovou que a integração iterativa de múltiplas ferramentas conseguiu reduzir a ocorrência de vulnerabilidades críticas geradas em até 96\%\ .

Com base nesses trabalhos, esta pesquisa implementou e avaliou um pipeline que combina LLMs, análise estática, testes automatizados e diferentes estratégias de contextualização e raciocínio guiado. A avaliação examinou o aceite das propostas individuais e o comportamento do projeto após a aplicação conjunta das alterações.